\origpage{103}{GRAPH MINORS. XI}
\noindent $c, b$-optimal. Since $c, b \in X$ and $X' \nsubseteq X$ (because $r' \subseteq X'$ and $r' \nsubseteq X$) this contradicts (8.9). Thus $r$ satisfies the theorem, as required. $\square$

Statement (8.12) turns out to be quite significant. For instance, the result of [5], that for planar graphs branch-width is computable in polynomial time, relies heavily on (8.12), and it will also have several applications in future papers of this series.

Another consequence of (8.9) is the following.

\stmt
(8.13) \textit{Let $a, b \in \Sigma$ with $\{a, b\} \cap U(\Gamma) \subseteq V(\Gamma)$; and let $X$ be a restraint with $a, b \in X$. Suppose that there is no restraint $Y$ with $X \subseteq Y$ and $l(Y) < l(X)$. Then $X$ includes an $a, b$-optimal restraint.}

\textit{Proof.} We have a slope $\operatorname{ins}$ in $\Gamma$ of order $\theta$. Let $\Gamma'$ be obtained from $\Gamma$ by subdividing each edge $e \in E(\Gamma)$ $2\theta - 1$ times if $e \subseteq X$, and $2\theta$ times otherwise. (Thus, $e$ corresponds to a path of $\Gamma'$ with $2\theta$ or $2\theta + 1$ edges.) Let $C$ be a circuit of $\Gamma$, and $C'$ the circuit of $\Gamma'$ with $U(C) = U(C')$. Then
\begin{equation*}
2\theta |E(C)| \le |E(C')| \le (2\theta + 1) |E(C)|
\end{equation*}
and so $C$ has length $<2\theta$ if and only if $C'$ has length $<4\theta^2$. If $C$ has length $<2\theta$, we define $\operatorname{ins}'(C') = \operatorname{ins}(C)$. Then $\operatorname{ins}'$ is a slope in $\Gamma'$ of order $4\theta^2$.

Define $\mathscr{R}$ to be the set of restraints of $\Gamma$ with respect to $\operatorname{ins}$, and $\mathscr{R}'$ to be the set of restraints of $\Gamma'$ with respect to $\operatorname{ins}'$ (replacing $2\theta$ by $4\theta^2$ in the definition of ``restraint''). Let us say that $X' \in \mathscr{R}'$ is \textit{smooth} if every vertex with valency $\le 1$ of $\Gamma' | bd(X')$ belongs to $V(\Gamma)$. Then a computation similar to that above yields that $\mathscr{R} \subseteq \mathscr{R}'$, and $\mathscr{R}$ contains precisely the smooth members of $\mathscr{R}'$.

Now $X \in \mathscr{R}$ and $l'(X) = 2\theta l(X)$, where $l'$ denotes the perimeter length of $X$ in $\Gamma'$. For any $Y \in \mathscr{R}'$ with $X \subseteq Y \ne X$, we claim that $l'(Y) > l'(X)$. For if $Y$ is smooth then $Y \in \mathscr{R}$; and since some edge of the perimeter of $Y$ is not included in $X$, and since from the hypothesis $l(Y) \ge l(X)$, it follows that $l'(Y) > 2\theta l(Y) \ge 2\theta l(X) = l'(X)$ as claimed. On the other hand, if $Y$ is not smooth then, since $X$ is smooth, there exists a smooth $Y' \in \mathscr{R}'$ with $X \subseteq Y' \subseteq Y$ and with $l'(Y') < l'(Y)$. But then $l'(Y') \ge l'(X)$, since either $Y' = X$ or $l'(Y') > l'(X)$. Thus $l'(Y) > l'(Y') \ge l'(X)$. This proves our claim. In particular, $X$ is extreme (in $\Gamma'$ with respect to $\operatorname{ins}'$). Let $X' \in \mathscr{R}'$ be an $a, b$-optimal restraint (in $\Gamma'$, with respect to $\operatorname{ins}'$). By (8.9), $X' \subseteq X$. Moreover, $X'$ is smooth, since $\{a, b\} \cap U(\Gamma) \subseteq V(\Gamma)$ and $X'$ is $a, b$-optimal; and hence $X' \in \mathscr{R}$. We claim that $X'$ is $a, b$-optimal in $\Gamma$ with respect to $\operatorname{ins}$. For suppose not, and let $Y \in \mathscr{R}$ with $a, b \in Y$ and $l(Y) < l(X')$. Then $Y \in \mathscr{R}'$, and
\begin{align*}
l'(Y) &\le (2\theta + 1) l(Y) \le (2\theta + 1)(l(X') - 1) \\
&= 2\theta l(X') + (l(X') - 2\theta - 1) < 2\theta l(X') \le l'(X')
\end{align*}
